\documentclass[10pt,twocolumn,letterpaper]{article}
\usepackage[margin=0.75in,columnsep=0.25in]{geometry}
\usepackage{amsmath,amssymb,amsthm,mathtools}
\usepackage[hidelinks]{hyperref}
\usepackage{microtype}
\newtheorem{theorem}{Theorem}
\newtheorem{lemma}{Lemma}
\theoremstyle{definition}
\newtheorem{definition}{Definition}
\newtheorem{example}{Example}
\theoremstyle{plain}
\setcounter{theorem}{5}
\setcounter{lemma}{6}
\setcounter{equation}{0}
\newcommand{\col}[1]{\begin{bmatrix}#1\end{bmatrix}}
\newcommand{\ip}[2]{\left\langle #1,#2\right\rangle}
\newcommand{\norm}[1]{\left\lVert #1\right\rVert}
\newcommand{\TT}{\mathcal T}
\newcommand{\opA}{\mathcal A}
\newcommand{\opB}{\mathcal B}
\newcommand{\opC}{\mathcal C}
\newcommand{\opD}{\mathcal D}
\newcommand{\opP}{\mathcal P}
\newcommand{\opK}{\mathcal K}
\allowdisplaybreaks
\title{\large Definitions, Theorem 6, and Lemma 7\\
\normalsize Nonzero plant initial conditions}
\author{}
\date{}
\begin{document}
\maketitle
\noindent
The following replacement text follows the notation, numbering, and proof
structure of the supplied ACC paper~\cite{channels}. Throughout, $G$ maps
an input to the output for the specified initial state $x(0)$; this initial
state is arbitrary and is never set to zero. Thus an initialized linear
plant generally defines an affine input--output operator. The multiplier
$\Psi$ has its standard zero initialization, and $\Delta$ has the
initialization specified by its input--output operator. No conclusion
about arbitrary hidden states of $\Delta$ is implied.

\section*{Notation and signal spaces}

For a Hilbert space $\mathcal H$, let $L_{\mathcal H}$ denote the space
of square-integrable functions $u:[0,\infty)\to\mathcal H$, with
\[
 \ip{u}{v}_{L_{\mathcal H}}
 =\int_0^\infty\ip{u(t)}{v(t)}_{\mathcal H}\,dt.
\]
Let $L_{e,\mathcal H}$ denote the extended space of functions for which
$\int_0^T\norm{u(t)}_{\mathcal H}^2dt$ is finite for every $T\geq0$.
Given $a<b$ and $n=(n_1,n_2)$, write
\[
 \begin{gathered}
 Z^n:=\mathbb R^{n_1}\times L_2^{n_2}[a,b],\\
 L^n_{e,[a,b]}:=L_{e,Z^n},\qquad L^n_Z:=L_{Z^n}.
 \end{gathered}
\]
The norm on $Z^n$ is induced by the Euclidean inner product and the
spatial $L_2$ inner product. The stacked signal $\col{u\\v}$ denotes
componentwise concatenation in $Z^{n+m}$. The subscript $L$ denotes the
appropriate time-signal Hilbert norm, and $Z$ the appropriate state or
spatial Hilbert norm. Define the time-truncation operator by
\[
 (P_\tau y)(t)=
 \begin{cases}y(t),&0\leq t\leq\tau,\\0,&t>\tau.\end{cases}
\]
The bounded PI operators $\Pi_4$ and their adjoints are as in the paper.
The storage operator $\opP$ is distinct from the truncation $P_\tau$.

\begin{definition}[Operators]
Let $\mathcal H,\mathcal G$ be Hilbert spaces. For each fixed admissible
initial condition $x(0)$, an operator
$G:L_{e,\mathcal H}\to L_{e,\mathcal G}$ is
\begin{enumerate}
\item \emph{Causal} if $P_\tau G=P_\tau GP_\tau$ for any $\tau\geq0$.
The same initial condition is used on both sides.
\item \emph{Bounded on $L_{\mathcal H}$ allowing for initial conditions}
if there exist $C_G,c_G\geq0$ such that, for all $v\in L_{\mathcal H}$
and all admissible $x(0)$,
\[
 \norm{Gv}_{L_{\mathcal G}}
 \leq C_G\norm{v}_{L_{\mathcal H}}+c_G\norm{x(0)}_Z.
\]
\item \emph{Bounded on $L_{e,\mathcal H}$ allowing for initial conditions}
if there exist $C_G,c_G\geq0$ such that, for all
$v\in L_{e,\mathcal H}$, all admissible $x(0)$, and all $\tau\geq0$,
\[
 \norm{P_\tau Gv}_{L_{\mathcal G}}
 \leq C_G\norm{P_\tau v}_{L_{\mathcal H}}+c_G\norm{x(0)}_Z.
\]
\item \emph{Incrementally $L_{e,\mathcal H}$-bounded} if there exists
$C\geq0$ such that, for all $v,u\in L_{e,\mathcal H}$,
\[
 \norm{P_\tau(Gv-Gu)}_{L_{\mathcal G}}
 \leq C\norm{P_\tau(v-u)}_{L_{\mathcal H}},\quad \tau\geq0.
\]
Here $Gv$ and $Gu$ have the same initial condition; the bound is uniform
over the admissible initial conditions under consideration.
\end{enumerate}
The constants in items 2 and 3 are independent of $x(0)$.
\end{definition}

Items 2 and 3 are the modifications to the paper's boundedness
definitions. For an operator without an initial-state contribution,
take $c_G=0$ to recover the original bounds. In particular, this
convention is used for $\Psi$. The notation $\mathcal L(\mathcal H,\mathcal G)$
continues to mean bounded linear operators; an initialized plant $G$
need not be linear as an input--output operator. For
$K\in\mathcal L(\mathcal H)$, define
\[
 (M_Kw)(t):=Kw(t).
\]

\section*{Feedback interconnections on Hilbert space}

\begin{definition}[Interconnection of $G$ and $\Delta$]
Given operators $G:L^n_{e,[a,b]}\to L^m_{e,[a,b]}$ and
$\Delta:L^m_{e,[a,b]}\to L^n_{e,[a,b]}$, an admissible initial condition
$x(0)$ for $G$, and signals $e\in L^n_{e,[a,b]}$ and
$f\in L^m_{e,[a,b]}$, we say that $u\in L^n_{e,[a,b]}$ and
$v\in L^m_{e,[a,b]}$ satisfy the interconnection defined by $[G,\Delta]$ if
\begin{equation}\label{eq:one}
 v=Gu+f,\qquad u=\Delta v+e.
\end{equation}
Here $G$ is evaluated with the specified initial condition $x(0)$.
\end{definition}

\begin{definition}[Well-posedness]
Given $G,\Delta$ as above, we say the interconnection defined by
$[G,\Delta]$ is well posed for all admissible plant initial conditions if,
for every such $x(0)$ and any $e\in L^n_{e,[a,b]}$,
$f\in L^m_{e,[a,b]}$, the following hold.
\begin{enumerate}
\item \emph{Existence and uniqueness:} There exist unique
$u\in L^n_{e,[a,b]}$, $v\in L^m_{e,[a,b]}$ satisfying
\eqref{eq:one} with that $x(0)$.
\item \emph{Causality:} If $u,v$ satisfy the interconnection for $e,f$,
and $\hat u,\hat v$ satisfy it for $P_\tau e,P_\tau f$, with the same
$x(0)$, then $P_\tau(u-\hat u)=0$ and $P_\tau(v-\hat v)=0$.
\end{enumerate}
\end{definition}

For a well-posed interconnection, retain the paper's notation
\[
 \col{u\\v}=F_{G,\Delta}\left(\col{e\\f}\right),
\]
with the dependence on the specified $x(0)$ left implicit in $G$.

\begin{definition}[Stability allowing for plant initial conditions]
We say the feedback system defined by $[G,\Delta]$ is stable allowing
for plant initial conditions if it is well posed in the sense of
Definition 3 and there exist constants $C,C_{\mathrm{init}}\geq0$ such
that, for any $e\in L^n_Z$, $f\in L^m_Z$, and every admissible $x(0)$,
the corresponding signals satisfy
\[
 \norm{\col{u\\v}}_L
 \leq C\norm{\col{e\\f}}_L+C_{\mathrm{init}}\norm{x(0)}_Z.
\]
Both constants are independent of $e,f,x(0)$. Stability is robust over
a class $\boldsymbol\Delta$ if well-posedness and this estimate hold
for every $\Delta\in\boldsymbol\Delta$ with common constants.
\end{definition}

This changes the paper's Definition 4 by retaining the possible output
caused by the initial state. It is an input--output definition; it does
not, by itself, assert convergence of hidden internal states.

\begin{definition}[Hard IQC]
We say the operator $\Delta:L^m_{e,[a,b]}\to L^n_{e,[a,b]}$ satisfies
the hard IQC defined by
$\Psi:L^{m+n}_{e,[a,b]}\to L^{m+n}_{e,[a,b]}$ and
$K\in\mathcal L(Z^{m+n})$ if, for any $T>0$ and all
$v\in L^m_{e,[a,b]}$,
\begin{equation}\label{eq:two}
 \ip{P_T\Psi\col{v\\\Delta v}}{
     P_TM_K\Psi\col{v\\\Delta v}}_L\geq0.
\end{equation}
The multiplier uses its specified zero initialization, and $\Delta$
uses the initialization specified by its input--output operator.
\end{definition}

Definition 5 is unchanged apart from making initialization and the
input dimension explicit. It does not require $x(0)=0$ for the plant.

\section*{The main theorem}

\begin{theorem}[IQC theorem with nonzero plant initial conditions]
If
\begin{enumerate}
\item $G:L^n_{e,[a,b]}\to L^m_{e,[a,b]}$ is causal and bounded on
$L^n_{e,[a,b]}$ allowing for initial conditions, in the sense of
Definition 1, with constants $C_G,c_G$;
\item $\Delta:L^m_{e,[a,b]}\to L^n_{e,[a,b]}$ and the interconnection
defined by $[G,\Delta]$ is well posed as in Definition 3;
\item $\Psi:L^{n+m}_{e,[a,b]}\to L^{n+m}_{e,[a,b]}$ is causal and
incrementally bounded, $\Psi(0)=0$, and $K\in\mathcal L(Z^{n+m})$;
\item $\Delta$ satisfies the hard IQC defined by $\Psi,K$;
\item there exist $\varepsilon>0$, $c_V\geq0$, and a nonnegative
function $V$ satisfying $V(x(0))\leq c_V\norm{x(0)}_Z^2$ such that,
for every admissible $x(0)$ and every $u\in L^n_{e,[a,b]}$,
\begin{equation}\label{eq:three}
\begin{aligned}
 &\ip{P_T\Psi\col{Gu\\u}}{
          P_TM_K\Psi\col{Gu\\u}}_L\\
 &\qquad\leq
 -\varepsilon\norm{P_Tu}_L^2+V(x(0)),
 \qquad T>0.
\end{aligned}
\end{equation}
\end{enumerate}
then the feedback system defined by $[G,\Delta]$ is stable in the sense
of Definition 4.
If these hypotheses hold with the same $G$ bounds, $\Psi,K$,
$\varepsilon,c_V$ for every $\Delta\in\boldsymbol\Delta$, then the
stability is robust over $\boldsymbol\Delta$ in the sense of Definition 4.
\end{theorem}

\begin{proof}
Define the shorthand notation $u_T:=P_Tu$ and
$\ip{u}{v}_T:=\ip{P_Tu}{P_Tv}_L$.
Clearly, by Cauchy--Schwarz,
$\ip{u}{v}_T\leq\norm{u_T}_L\norm{v_T}_L$.

Now, because the interconnection defined by $[G,\Delta]$ is well posed,
it suffices to establish the claimed bound uniformly over all finite
horizons. For $e\in L^n_Z$, $f\in L^m_Z$, suppose that $u,v$ satisfy
the interconnection defined by $G,\Delta$. Then
\[
 v=Gu+f,\qquad u=\Delta v+e.
\]
By the hard IQC and \eqref{eq:three},
\begin{align*}
 \varepsilon\norm{u_T}_L^2
 &\leq V(x(0))
  +\ip{\Psi\col{v\\\Delta v}}{M_K\Psi\col{v\\\Delta v}}_T\\
 &\quad-\ip{\Psi\col{Gu\\u}}{M_K\Psi\col{Gu\\u}}_T.
\end{align*}

Now, because $\Psi$ and hence $M_K\Psi$ are incrementally bounded,
with bounds $C_\Psi$ and
$C_{M_K\Psi}=C_\Psi\norm{K}$, respectively, and
$\Psi(0)=0$, expanding the difference of the quadratic forms and
applying Cauchy--Schwarz gives
\begin{align*}
 \varepsilon\norm{u_T}_L^2
 &\leq V(x(0))+C_\Psi C_{M_K\Psi}
 \left(\norm{\col{v-Gu\\\Delta v-u}_T}_L^2\right.\\
 &\left.\qquad
 +2\norm{\col{Gu\\u}_T}_L
      \norm{\col{v-Gu\\\Delta v-u}_T}_L\right).
\end{align*}
Since $v-Gu=f$, $\Delta v-u=-e$, and
\[
 \norm{\col{Gu\\u}_T}_L
 \leq (C_G+1)\norm{u_T}_L+c_G\norm{x(0)}_Z,
\]
we conclude that
\begin{align*}
 \varepsilon\norm{u_T}_L^2
 &\leq V(x(0))
 +C_\Psi C_{M_K\Psi}\norm{\col{f\\e}_T}_L^2\\
 &\quad+2C_\Psi C_{M_K\Psi}(C_G+1)
       \norm{u_T}_L\norm{\col{f\\e}_T}_L\\
 &\quad+2C_\Psi C_{M_K\Psi}c_G\norm{x(0)}_Z
                         \norm{\col{f\\e}_T}_L.
\end{align*}
Applying Young's inequality to the last two terms yields
\begin{align*}
 \frac{\varepsilon}{2}\norm{u_T}_L^2
 &\leq V(x(0))
      +C_\Psi C_{M_K\Psi}c_G^2\norm{x(0)}_Z^2\\
 &\quad+2C_\Psi C_{M_K\Psi}\norm{\col{f\\e}_T}_L^2\\
 &\quad+\frac{2(C_\Psi C_{M_K\Psi})^2(C_G+1)^2}{\varepsilon}
                  \norm{\col{f\\e}_T}_L^2.
\end{align*}
Using $V(x(0))\leq c_V\norm{x(0)}_Z^2$, define
\begin{align*}
 \gamma_x^2&=\frac{2(c_V+C_\Psi C_{M_K\Psi}c_G^2)}
                  {\varepsilon},\\
 \gamma_d^2&=\frac{4C_\Psi C_{M_K\Psi}}{\varepsilon}
 +\frac{4(C_\Psi C_{M_K\Psi})^2(C_G+1)^2}{\varepsilon^2}.
\end{align*}
Thus
\[
 \norm{u_T}_L\leq
 \gamma_d\norm{\col{f\\e}_T}_L+\gamma_x\norm{x(0)}_Z.
\]
Finally, since $v=Gu+f$,
\begin{align*}
 \norm{\col{u\\v}_T}_L
 &\leq (C_G+1)\norm{u_T}_L+
                 c_G\norm{x(0)}_Z+\norm{f_T}_L\\
 &\leq \bigl((C_G+1)\gamma_d+1\bigr)
                         \norm{\col{f\\e}_T}_L\\
 &\quad+\bigl((C_G+1)\gamma_x+c_G\bigr)\norm{x(0)}_Z.
\end{align*}
The constants are independent of $T$. Letting $T\to\infty$ proves the
claim. No step sets $x(0)=0$.
\end{proof}

\noindent\emph{Remark:}
The bound on $G$ replaces the original zero-state boundedness assumption.
For a linear plant it requires a uniformly finite-energy initial output
in addition to the bounded zero-state input--output map. It does not
follow from zero-state boundedness alone. The conclusion is initialized
input--output stability, not an assertion of internal state convergence.

\medskip
\noindent
As in the paper, consider a PIE realization
\begin{equation}\label{eq:four}
 \begin{aligned}
 \TT\dot x(t)&=\opA x(t)+\opB u(t),\\
 y(t)&=\opC x(t)+\opD u(t).
 \end{aligned}
\end{equation}
The initial condition is now any admissible $x(0)$.

\begin{lemma}[Sufficient version of KYP lemma with nonzero initial conditions]
If
\begin{enumerate}
\item $G:L^n_{e,[a,b]}\to L^m_{e,[a,b]}$ is causal and bounded allowing
for initial conditions in the sense of Definition 1;
\item there exist $\TT,\opA,\opB,\opC,\opD\in\Pi_4$ such that, for any
$u\in L^n_{e,[a,b]}$, $y=Gu$ implies that $y$ satisfies
\eqref{eq:four} for a trajectory $x$ with the specified admissible
$x(0)$, and the integrated dissipation calculation below is valid
for these trajectories;
\item $K\in\Pi_4$ and there exist $\varepsilon>0$ and
$\opP\in\Pi_4$ such that $\opP\succeq0$ and
\begin{equation}\label{eq:five}
\begin{aligned}
 &\col{\TT^*\opP\opA+\opA^*\opP\TT&
                       \TT^*\opP\opB\\
        \opB^*\opP\TT&\varepsilon I}\\
 &\quad+\col{\opC^*\\\opD^*}
             K\begin{bmatrix}\opC&\opD\end{bmatrix}
 \preceq0.
\end{aligned}
\end{equation}
\end{enumerate}
then, for any admissible $x(0)$ and any
$u\in L^n_{e,[a,b]}$,
\begin{equation*}
 \begin{aligned}
 &\ip{P_TGu}{P_TM_KGu}_L\\
 &\quad\leq-\varepsilon\norm{P_Tu}_L^2
       +\ip{\TT x(0)}{\opP\TT x(0)}_Z,
 \end{aligned}
\end{equation*}
for all $T>0$.
\end{lemma}

\begin{proof}
Define $V(x)=\ip{\TT x}{\opP\TT x}$. Since $\opP\succeq0$ and
$\TT,\opP\in\Pi_4$ are bounded, Cauchy--Schwarz and the operator
norm give
\[
 0\leq V(x)=\ip{x}{\TT^*\opP\TT x}
 \leq\norm{\TT^*\opP\TT}\norm{x}_Z^2.
\]
Suppose that
$u\in L^n_{e,[a,b]}$ and $y(t)=(Gu)(t)$ for a trajectory
$x$ satisfying \eqref{eq:four} with arbitrary admissible $x(0)$.
By inequality \eqref{eq:five}, multiplying on both sides by
$\col{x(t)\\u(t)}$ in the quadratic form gives
\begin{align*}
 &\ip{\TT\dot x(t)}{\opP\TT x(t)}
   +\ip{\TT x(t)}{\opP\TT\dot x(t)}\\
 &\quad+\varepsilon\norm{u(t)}_Z^2+\ip{y(t)}{Ky(t)}_Z\\
 &=\dot V(x(t))+\varepsilon\norm{u(t)}_Z^2
                         +\ip{y(t)}{Ky(t)}_Z
 \leq0.
\end{align*}
Integrating in time, we obtain
\begin{align*}
 &\varepsilon\norm{P_Tu}_L^2
    +\int_0^T\ip{y(t)}{Ky(t)}_Z\,dt\\
 &\qquad\leq -V(x(T))+V(x(0))\\
 &\qquad\leq V(x(0)),
\end{align*}
since $V(x(T))\geq0$. We conclude that
\begin{align*}
 &\ip{P_TGu}{P_TM_KGu}_L\\
 &=\int_0^T\ip{y(t)}{Ky(t)}_Z\,dt\\
 &\leq-\varepsilon\norm{P_Tu}_L^2+V(x(0))\\
 &=-\varepsilon\norm{P_Tu}_L^2
       +\ip{\TT x(0)}{\opP\TT x(0)}_Z.
\end{align*}
The initial term is retained throughout.
\end{proof}

\noindent\emph{Remarks:}
The cross blocks in \eqref{eq:five} include
$\TT^*\opP\opB$ and $\opB^*\opP\TT$, as required by the storage
$V(x)=\ip{\TT x}{\opP\TT x}$ and the stated dynamics.
The supplied PDF prints $\opP\opB$ and $\opB^*\opP$ instead.
For mild solutions, the integrated inequality must be justified by
regularity or extension from regular solutions.

Applying this lemma to the augmented plant--multiplier realization, as in
Corollary 8 of the original paper, supplies \eqref{eq:three}, with the
initial-storage allowance evaluated at
$\col{x(0)\\0}$. The initialized bound on the nominal $G$ is still
a separate hypothesis; the lemma does not establish it.

One way to compute an admissible $c_V$ is to impose
\[
 0\preceq\TT^*\opP\TT\preceq c_V I
\]
together with \eqref{eq:five}. Since $\TT$ is fixed, the upper bound
is linear in $\opP,c_V$. Alternatively take
$c_V=\norm{\TT}^2\norm{\opP}$. No coercive lower bound on the storage
is needed for the initialized input--output conclusion of Theorem 6.

\section*{Examples verifying the conditions}

\begin{example}[Verifying initialized boundedness by a gain LPI]
For a PIE realization \eqref{eq:four}, suppose one can find
$\mathcal R\in\Pi_4$, $\mathcal R\succeq0$, and constants
$\gamma,r_0\geq0$ satisfying
\begin{gather*}
 \col{\TT^*\mathcal R\opA+\opA^*\mathcal R\TT&
          \TT^*\mathcal R\opB\\
       \opB^*\mathcal R\TT&-\gamma^2I}
 +\col{\opC^*\\\opD^*}
       \begin{bmatrix}\opC&\opD\end{bmatrix}\preceq0,\\
 0\preceq\TT^*\mathcal R\TT\preceq r_0 I.
\end{gather*}
Multiplication by $\col{x(t)\\u(t)}$ and integration, justified for
the chosen realization, give
\[
 \norm{P_TGu}_L^2
 \leq\gamma^2\norm{P_Tu}_L^2+r_0\norm{x(0)}_Z^2.
\]
Consequently Definition 1(3) holds with
$C_G=\gamma$, $c_G=\sqrt{r_0}$. These are separate sufficient gain
conditions for Theorem 6's first assumption. The operator
$\mathcal R$ need not equal the IQC storage $\opP$. For fixed
$\gamma$, these are LPIs in $\mathcal R,r_0$; one can also search
over $\gamma^2$ as a scalar variable.
\end{example}

\begin{example}[The initialized bound for a stable nominal realization]
Consider a linear realization on a Hilbert state space, with bounded
input and output operators $B,C,D$. Suppose its state-transition
semigroup satisfies
\[
 \norm{S(t)}\leq M e^{-\alpha t},\qquad M,\alpha>0.
\]
For any admissible initial state, its output is
\[
 (Gu)(t)=CS(t)x(0)+Du(t)
             +\int_0^t CS(t-s)Bu(s)\,ds.
\]
The initial-output term satisfies
\[
 \norm{CS(\cdot)x(0)}_L
 \leq\frac{M\norm{C}}{\sqrt{2\alpha}}\norm{x(0)}_Z.
\]
The convolution kernel has $L_1$ norm at most
$M\norm{C}\norm{B}/\alpha$. Young's convolution inequality, applied
to the truncated input, therefore verifies Definition 1(3) with
\begin{align*}
 C_G&=\norm{D}+\frac{M\norm{C}\norm{B}}{\alpha},\\
 c_G&=\frac{M\norm{C}}{\sqrt{2\alpha}}.
\end{align*}
This is a sufficient way to verify the added assumption, not an
assertion that exponential nominal stability is necessary. If the
realization has an unbounded input or observation operator, appropriate
admissibility estimates must replace these bounded-operator formulas.
\end{example}

\begin{example}[A complete scalar certificate with arbitrary initial state]
Consider
\[
 \dot x(t)=-2x(t)+u(t),\qquad y(t)=x(t),\qquad x(0)\in\mathbb R.
\]
Example 2 gives $C_G=1/2$ and $c_G=1/2$. Let
$\boldsymbol\Delta$ consist of the memoryless operators
$(\Delta v)(t)=\phi(v(t))$, where $\phi(0)=0$ and $\phi$ is globally
Lipschitz with constant at most one. Then $|\phi(v)|\leq|v|$, so every
member satisfies the same hard IQC with
\[
 \Psi=I,\qquad K=\col{1&0\\0&-1}.
\]
Indeed,
\[
 \int_0^T\bigl(|v(t)|^2-|\phi(v(t))|^2\bigr)dt\geq0.
\]
Here $\Psi$ is causal, $\Psi(0)=0$, and its incremental bound is one.

For arbitrary $e,f\in L_{2e}$, the feedback equations give
\[
 \dot x(t)=-2x(t)+\phi(x(t)+f(t))+e(t).
\]
Its right-hand side is globally Lipschitz in $x$, uniformly in time,
and has locally integrable forcing. Consequently it has a unique
global absolutely continuous solution for every $x(0)$; causality
follows from uniqueness. This verifies Definition 3 without requiring
an algebraic-loop assumption, because the plant has no feedthrough.

Apply Lemma 7 to the graph output $\col{y\\u}$, as in the paper's
augmented construction. Its output operators are
$\opC=\col{1\\0}$ and $\opD=\col{0\\1}$, while
$\TT=1$, $\opA=-2$, and $\opB=1$. Choose
\[
 \opP=2,\qquad \varepsilon=\tfrac14.
\]
The left-hand side of \eqref{eq:five} becomes
\[
 \col{-8&2\\2&1/4}+\col{1&0\\0&-1}
 =\col{-7&2\\2&-3/4}\prec0.
\]
The first diagonal entry is negative and the determinant is
$21/4-4=5/4>0$, verifying strict negative definiteness. The storage is
\[
 V(x)=2x^2,\qquad c_V=2.
\]
Lemma 7 consequently supplies precisely the initialized separation
condition required in Theorem 6:
\[
 \int_0^T\bigl(|(Gu)(t)|^2-|u(t)|^2\bigr)dt
 \leq-\tfrac14\int_0^T|u(t)|^2dt+2|x(0)|^2.
\]
All constants are common to the uncertainty class. Theorem 6 therefore
gives robust input--output stability allowing for every $x(0)$.
The time-domain calculations in Lemma 7 are valid almost everywhere
for these absolutely continuous trajectories and can be integrated.
\end{example}

\begin{example}[An infinite-dimensional initial-output bound]
On the Hilbert space $L_2(0,1)$ consider the heat equation
\begin{align*}
 x_t(s,t)&=x_{ss}(s,t)+u(s,t),\\
 x(0,t)&=x(1,t)=0,\qquad y(s,t)=x(s,t).
\end{align*}
The initial profile $x(\cdot,0)$ may be any element of $L_2(0,1)$.
The Dirichlet heat semigroup has eigenvalues $-j^2\pi^2$,
$j=1,2,\ldots$, and satisfies $\norm{S(t)}\leq e^{-\pi^2t}$.
Taking $B=C=I$, $D=0$ in Example 2 gives
\[
 C_G=\frac{1}{\pi^2},\qquad
 c_G=\frac{1}{\sqrt{2}\,\pi}.
\]
Thus the initialized boundedness assumption is meaningful for arbitrary
nonzero infinite-dimensional initial profiles, not just for ODE states.

For a pointwise nonlinearity $\phi$ as in Example 3, the induced operator
on $L_2(0,1)$ is globally Lipschitz with constant one. The resulting
semilinear heat equation with $e,f\in L_{2e}$ has unique global mild
solutions. This verifies the corresponding well-posedness requirement.
For regular initial profiles and inputs, energy differentiation is
available directly. To use a PIE storage estimate for all mild
solutions, first prove it for regular data and pass to the limit using
density, continuous dependence, and boundedness of the storage
operators. A particular PIE realization and its LPI must still be
verified; the semigroup bound alone does not verify \eqref{eq:five}.
The initial-state norm and admissible profiles must match that PIE
realization: a general physical $L_2$ profile is not automatically an
admissible initial fundamental state if that state contains derivatives.
\end{example}

\begin{example}[A dynamic multiplier satisfying the IQC conditions]
Let $(\Delta v)(t)=d\,v(t)$ for an unknown constant $d\in[-1,1]$.
Choose the zero-initialized stable filter $H$ with transfer function
$H(s)=1/(s+1)$, acting componentwise if the signals are Hilbert-valued,
and take
\[
 \Psi=\col{H&0\\0&H},\qquad
 K=\col{I&0\\0&-I}.
\]
The filter is causal, bounded, incrementally bounded, and maps zero to
zero. Its $L_2$ gain is at most one. Since $d$ is constant and $H$ is
linear, $H(dv)=dHv$. Therefore
\begin{align*}
 &\ip{P_T\Psi\col{v\\\Delta v}}{
          P_TM_K\Psi\col{v\\\Delta v}}_L\\
 &\qquad=(1-d^2)\int_0^T\norm{(Hv)(t)}_Z^2dt\geq0.
\end{align*}
This verifies the hard IQC with the same multiplier for every allowed
$d$. The argument uses a constant scalar uncertainty; it does not
establish the same filtered IQC for a general time-varying $d(t)$.
The plant still needs the initialized bound, well-posedness, and the
separation certificate in Theorem 6.
\end{example}

\begin{thebibliography}{1}
\bibitem{channels}
A. Talitckii, M. M. Peet, and P. Seiler,
\emph{Integral Quadratic Constraints with Infinite-Dimensional Channels},
2023 American Control Conference, pp. 1576--1583, 2023,
doi: 10.23919/ACC55779.2023.10156335.
\end{thebibliography}
\end{document}

